\section{Conclusion}
\label{sec:conclusion}

We asked whether SOCIAL's synthetic small-world communication graph should be replaced by a chemically meaningful MOF similarity network when the optimiser is used to select DFT evaluations under a budget. Over 8{,}697 QMOF MOFs, four band-gap objectives, four budgets and 30 seeds, the answer is no, at least not detectably. The pre-specified pilot test of H1 failed. The 30-seed comparison found at most small, objective-dependent advantages over a degree-preserving random topology ($\delta\leq0.22$, power $\leq0.64$, not significant after correction), and a Watts--Strogatz graph did better on one objective. High-betweenness neighbours did not lead to new chemical families more often than expected (H3). \gs{} beat random search only on the minimum-gap objective, and that advantage came from its chemistry-aware embedding, not from the graph. Surrogate-based Thompson sampling and snap-based PSO, DE and GA outperformed it.

Along the way, we found that the published MOFGalaxyNet statistics correspond to a threshold near 0.7 and to a different recipe from the one described, and that MOFGalaxyNet graphs are strongly band-gap homophilous, largely because of building-block duplicates. Future work could couple agent movement directly to graph edges, combine graph communication with a surrogate model, and test topologies on objectives whose hits are concentrated in fewer chemical families. The complete gated pipeline, including all negative results, is released for reuse.
